\documentclass{amsart}
% For dealing with references we use the comment environment
\usepackage{verbatim}
\newenvironment{reference}{\comment}{\endcomment}
%\newenvironment{reference}{}{}
\newenvironment{slogan}{\comment}{\endcomment}
\newenvironment{history}{\comment}{\endcomment}

% For commutative diagrams we use Xy-pic
\usepackage[all]{xy}

% We use 2cell for 2-commutative diagrams.
\xyoption{2cell}
\UseAllTwocells

% We use multicol for the list of chapters between chapters
\usepackage{multicol}

% This is generally recommended for better output
\usepackage{lmodern}
\usepackage[T1]{fontenc}

% For cross-file-references

% Package for hypertext links:
\usepackage{hyperref}

% For any local file, say "hello.tex" you want to link to please

% Theorem environments.
%
\theoremstyle{plain}
\newtheorem{theorem}[subsection]{Theorem}
\newtheorem{proposition}[subsection]{Proposition}
\newtheorem{lemma}[subsection]{Lemma}

\theoremstyle{definition}
\newtheorem{definition}[subsection]{Definition}
\newtheorem{example}[subsection]{Example}
\newtheorem{exercise}[subsection]{Exercise}
\newtheorem{situation}[subsection]{Situation}

\theoremstyle{remark}
\newtheorem{remark}[subsection]{Remark}
\newtheorem{remarks}[subsection]{Remarks}

\numberwithin{equation}{subsection}

% Macros
%
\def\lim{\mathop{\mathrm{lim}}\nolimits}
\def\colim{\mathop{\mathrm{colim}}\nolimits}
\def\Spec{\mathop{\mathrm{Spec}}}
\def\Hom{\mathop{\mathrm{Hom}}\nolimits}
\def\Ext{\mathop{\mathrm{Ext}}\nolimits}
\def\SheafHom{\mathop{\mathcal{H}\!\mathit{om}}\nolimits}
\def\SheafExt{\mathop{\mathcal{E}\!\mathit{xt}}\nolimits}
\def\Sch{\mathit{Sch}}
\def\Mor{\mathop{\mathrm{Mor}}\nolimits}
\def\Ob{\mathop{\mathrm{Ob}}\nolimits}
\def\Sh{\mathop{\mathit{Sh}}\nolimits}
\def\NL{\mathop{N\!L}\nolimits}
\def\CH{\mathop{\mathrm{CH}}\nolimits}
\def\proetale{{pro\text{-}\acute{e}tale}}
\def\etale{{\acute{e}tale}}
\def\QCoh{\mathit{QCoh}}
\def\Ker{\mathop{\mathrm{Ker}}}
\def\Im{\mathop{\mathrm{Im}}}
\def\Coker{\mathop{\mathrm{Coker}}}
\def\Coim{\mathop{\mathrm{Coim}}}

% Boxtimes
%
\DeclareMathSymbol{\boxtimes}{\mathbin}{AMSa}{"02}

%
% Macros for moduli stacks/spaces
%
\def\QCohstack{\mathcal{QC}\!\mathit{oh}}
\def\Cohstack{\mathcal{C}\!\mathit{oh}}
\def\Spacesstack{\mathcal{S}\!\mathit{paces}}
\def\Quotfunctor{\mathrm{Quot}}
\def\Hilbfunctor{\mathrm{Hilb}}
\def\Curvesstack{\mathcal{C}\!\mathit{urves}}
\def\Polarizedstack{\mathcal{P}\!\mathit{olarized}}
\def\Complexesstack{\mathcal{C}\!\mathit{omplexes}}
% \Pic is the operator that assigns to X its picard group, usage \Pic(X)
% \Picardstack_{X/B} denotes the Picard stack of X over B
% \Picardfunctor_{X/B} denotes the Picard functor of X over B
\def\Pic{\mathop{\mathrm{Pic}}\nolimits}
\def\Picardstack{\mathcal{P}\!\mathit{ic}}
\def\Picardfunctor{\mathrm{Pic}}
\def\Deformationcategory{\mathcal{D}\!\mathit{ef}}

\setlength{\emergencystretch}{3em}
\begin{document}
\title[Stacks Project, Tag 0ESZ]{453 --- Chern invariants of locally bounded vector-bundle complexes depend only on their derived object}
\maketitle
\noindent Copyright (C) 2005--2025 Johan de Jong. Stacks Project authors. Source: \href{https://stacks.math.columbia.edu/tag/0ESZ}{Tag 0ESZ}.

\noindent Editorial difficulty note (not author text). Difficulty H2: Two representatives need not be termwise comparable globally, so the proof modifies each integral test scheme until the cohomology has Tor dimension at most one. It then compares the resulting vector-bundle resolutions through a Schanuel-type fibre-product construction, using bivariant detection to recover global independence..

\noindent Editorial provenance note (not author text). History: excerpt from revision \texttt{a0444\allowbreak 6e57e\allowbreak c1fbc\allowbreak 252a8\allowbreak 71afc\allowbreak ec775\allowbreak 2fb28\allowbreak 07b14}, chow.tex, lines 9090--9210. Reference presentation was adapted; original-author context and core support blocks were retained or added. The mathematical wording is unchanged. Licensed under GNU FDL 1.2 or later; no invariant sections or cover texts. See ../licenses/COPYING. Context blocks reproduce author definitions and supporting results; they are not additional counted problems.

\begin{definition}
\label{definition-bivariant-class}
\begin{reference}
Similar to \href{https://stacks.math.columbia.edu/bibliography}{Bibliography: F (Definition 17.1)}
\end{reference}
Let $(S, \delta)$ be as in Situation \ref{situation-setup}.
Let $f : X \to Y$ be a morphism of schemes locally of finite type over $S$.
Let $p \in \mathbf{Z}$.
A {\it bivariant class $c$ of degree $p$ for $f$} is given by a rule
which assigns to every locally of finite type morphism $Y' \to Y$
and every $k$ a map
$$
c \cap - : \CH_k(Y') \longrightarrow \CH_{k - p}(X')
$$
where $X' = Y' \times_Y X$, satisfying the following conditions
\begin{enumerate}
\item if $Y'' \to Y'$ is a proper, then
$c \cap (Y'' \to Y')_*\alpha'' = (X'' \to X')_*(c \cap \alpha'')$
for all $\alpha''$ on $Y''$ where $X'' = Y'' \times_Y X$,
\item if $Y'' \to Y'$ is flat locally of finite type of
fixed relative dimension, then
$c \cap (Y'' \to Y')^*\alpha' = (X'' \to X')^*(c \cap \alpha')$
for all $\alpha'$ on $Y'$, and
\item if $(\mathcal{L}', s', i' : D' \to Y')$ is as in
Definition \ref{definition-gysin-homomorphism}
with pullback $(\mathcal{N}', t', j' : E' \to X')$ to $X'$,
then we have $c \cap (i')^*\alpha' = (j')^*(c \cap \alpha')$
for all $\alpha'$ on $Y'$.
\end{enumerate}
The collection of all bivariant classes of degree $p$ for $f$ is
denoted $A^p(X \to Y)$.
\end{definition}

\begin{situation}
\label{situation-setup}
Here $S$ is a locally Noetherian, and universally catenary scheme.
Moreover, we assume $S$ is endowed with a dimension function
$\delta : S \longrightarrow \mathbf{Z}$.
\end{situation}

\begin{definition}
\label{definition-gysin-homomorphism}
Let $(S, \delta)$ be as in Situation \ref{situation-setup}.
Let $X$ be locally of finite type over $S$.
Let $(\mathcal{L}, s)$ be a pair consisting of an invertible
sheaf and a global section $s \in \Gamma(X, \mathcal{L})$.
Let $D = Z(s)$ be the zero scheme of $s$, and
denote $i : D \to X$ the closed immersion.
We define, for every integer $k$, a {\it Gysin homomorphism}
$$
i^* : Z_{k + 1}(X) \to \CH_k(D).
$$
by the following rules:
\begin{enumerate}
\item Given an integral closed subscheme $W \subset X$ with
$\dim_\delta(W) = k + 1$ we define
\begin{enumerate}
\item if $W \not \subset D$, then $i^*[W] = [D \cap W]_k$ as a
$k$-cycle on $D$, and
\item if $W \subset D$, then
$i^*[W] = i'_*(c_1(\mathcal{L}|_W) \cap [W])$,
where $i' : W \to D$ is the induced closed immersion.
\end{enumerate}
\item For a general $(k + 1)$-cycle $\alpha = \sum n_j[W_j]$
we set
$$
i^*\alpha = \sum n_j i^*[W_j]
$$
\item If $D$ is an effective Cartier divisor, then we denote
$D \cdot \alpha = i_*i^*\alpha$ the pushforward of the class $i^*\alpha$
to a class on $X$.
\end{enumerate}
\end{definition}

\begin{lemma}
\label{lemma-pre-derived-chern-class}
Let $(S, \delta)$ be as in Situation \ref{situation-setup}.
Let $X$ be locally of finite type over $S$. Let $E \in D(\mathcal{O}_X)$
be an object such that there exists a locally bounded complex
$\mathcal{E}^\bullet$ of finite locally free $\mathcal{O}_X$-modules
representing $E$. Then a slight generalization of the above constructions
$$
c(\mathcal{E}^\bullet) \in \prod\nolimits_{p \geq 0} A^p(X),\quad
ch(\mathcal{E}^\bullet) \in
\prod\nolimits_{p \geq 0} A^p(X) \otimes \mathbf{Q},\quad
P_p(\mathcal{E}^\bullet) \in A^p(X)
$$
are independent of the choice of the complex $\mathcal{E}^\bullet$.
\end{lemma}

\begin{proof}
We prove this for the total Chern class; the other two cases follow
by the same arguments using
Lemma \href{https://stacks.math.columbia.edu/tag/0F9C}{lemma chern character additive (Tag 0F9C)}
instead of
Lemma \href{https://stacks.math.columbia.edu/tag/02UI}{lemma additivity chern classes (Tag 02UI)}.

\medskip\noindent
As in Remark \href{https://stacks.math.columbia.edu/tag/0ESW}{remark extend to finite locally free (Tag 0ESW)} in order
to define the total chern class $c(\mathcal{E}^\bullet)$
we decompose $X$ into open and closed subschemes
$$
X = \coprod\nolimits_{i \in I} X_i
$$
such that the rank $\mathcal{E}^n$ is constant on $X_i$ for
all $n$ and $i$. (Since these ranks are locally constant functions
on $X$ we can do this.) Since $\mathcal{E}^\bullet$ is locally
bounded, we see that only a finite number of the sheaves
$\mathcal{E}^n|_{X_i}$ are nonzero for a fixed $i$. Hence we
can define
$$
c(\mathcal{E}^\bullet|_{X_i}) =
\prod\nolimits_n c(\mathcal{E}^n|_{X_i})^{(-1)^n}
\in \prod\nolimits_{p \geq 0} A^p(X_i)
$$
as above. By Lemma \href{https://stacks.math.columbia.edu/tag/0FDZ}{lemma disjoint decomposition bivariant (Tag 0FDZ)}
we have $A^p(X) = \prod_i A^p(X_i)$. Hence for each $p \in \mathbf{Z}$
we have a unique element $c_p(\mathcal{E}^\bullet) \in A^p(X)$ restricting
to $c_p(\mathcal{E}^\bullet|_{X_i})$ on $X_i$ for all $i$.

\medskip\noindent
Suppose we have a second locally bounded complex
$\mathcal{F}^\bullet$ of finite locally free $\mathcal{O}_X$-modules
representing $E$.
Let $g : Y \to X$ be a morphism locally of finite type with $Y$ integral.
By Lemma \href{https://stacks.math.columbia.edu/tag/02UC}{lemma bivariant zero (Tag 02UC)} it suffices to show that
with $c(g^*\mathcal{E}^\bullet) \cap [Y]$ is the same as
$c(g^*\mathcal{F}^\bullet) \cap [Y]$ and it even suffices to prove
this after replacing $Y$ by an integral scheme proper and birational
over $Y$. Then first we conclude that $g^*\mathcal{E}^\bullet$
and $g^*\mathcal{F}^\bullet$ are bounded complexes of finite locally
free $\mathcal{O}_Y$-modules of constant rank. Next, by
More on Flatness, Lemma \href{https://stacks.math.columbia.edu/tag/0ESR}{lemma blowup complex integral (Tag 0ESR)}
we may assume that $H^i(Lg^*E)$ is perfect of tor dimension $\leq 1$
for all $i \in \mathbf{Z}$.
This reduces us to the case discussed in the next paragraph.

\medskip\noindent
Assume $X$ is integral, $\mathcal{E}^\bullet$ and $\mathcal{F}^\bullet$
are bounded complexes of finite locally free modules of constant rank, and
$H^i(E)$ is a perfect $\mathcal{O}_X$-module of tor dimension $\leq 1$
for all $i \in \mathbf{Z}$. We have to
show that $c(\mathcal{E}^\bullet) \cap [X]$ is the same as
$c(\mathcal{F}^\bullet) \cap [X]$. Denote
$d_\mathcal{E}^i : \mathcal{E}^i \to \mathcal{E}^{i + 1}$ and
$d_\mathcal{F}^i : \mathcal{F}^i \to \mathcal{F}^{i + 1}$
the differentials of our complexes. By
More on Flatness, Remark \href{https://stacks.math.columbia.edu/tag/0F8J}{remark when you have a complex (Tag 0F8J)}
we know that $\Im(d_\mathcal{E}^i)$, $\Ker(d_\mathcal{E}^i)$,
$\Im(d_\mathcal{F}^i)$, and $\Ker(d_\mathcal{F}^i)$
are finite locally free $\mathcal{O}_X$-modules for all $i$.
By additivity (Lemma \href{https://stacks.math.columbia.edu/tag/02UI}{lemma additivity chern classes (Tag 02UI)}) we see that
$$
c(\mathcal{E}^\bullet) = \prod\nolimits_i
c(\Ker(d_\mathcal{E}^i))^{(-1)^i} c(\Im(d_\mathcal{E}^i))^{(-1)^i}
$$
and similarly for $\mathcal{F}^\bullet$. Since we have the
short exact sequences
$$
0 \to \Im(d_\mathcal{E}^i) \to \Ker(d_\mathcal{E}^i) \to H^i(E) \to 0
\quad\text{and}\quad
0 \to \Im(d_\mathcal{F}^i) \to \Ker(d_\mathcal{F}^i) \to H^i(E) \to 0
$$
we reduce to the problem stated and solved in the next paragraph.

\medskip\noindent
Assume $X$ is integral and we have two short exact sequences
$$
0 \to \mathcal{E}' \to \mathcal{E} \to \mathcal{Q} \to 0
\quad\text{and}\quad
0 \to \mathcal{F}' \to \mathcal{F} \to \mathcal{Q} \to 0
$$
with $\mathcal{E}$, $\mathcal{E}'$, $\mathcal{F}$, $\mathcal{F}'$
finite locally free. Problem: show that
$c(\mathcal{E})c(\mathcal{E}')^{-1} \cap [X] =
c(\mathcal{F})c(\mathcal{F}')^{-1} \cap [X]$.
To do this, consider the short exact sequence
$$
0 \to \mathcal{G} \to \mathcal{E} \oplus \mathcal{F} \to \mathcal{Q} \to 0
$$
defining $\mathcal{G}$. Since $\mathcal{Q}$ has tor dimension $\leq 1$
we see that $\mathcal{G}$ is finite locally free. A diagram chase
shows that the kernel of the surjection $\mathcal{G} \to \mathcal{F}$
maps isomorphically to $\mathcal{E}'$ in $\mathcal{E}$ and
the kernel of the surjection $\mathcal{G} \to \mathcal{E}$ maps
isomorphically to $\mathcal{F}'$ in $\mathcal{F}$. (Working affine
locally this follows from or is equivalent to Schanuel's lemma, see
Algebra, Lemma \href{https://stacks.math.columbia.edu/tag/00O3}{lemma Schanuel (Tag 00O3)}.)
We conclude that
$$
c(\mathcal{E})c(\mathcal{F}') = c(\mathcal{G}) =
c(\mathcal{F})c(\mathcal{E}')
$$
as desired.
\end{proof}
\end{document}
